% !TeX root = ../main.tex
\section{Metatheory}\label{sec:meta}

In this section, we provide the denotation of types ($\denot{\tau}$)
and type contexts ($\denot{\Gamma}$), as well as the soundness theorem
of our type system.  Building such a metatheory is challenging because
it must provide semantics for both demonic and angelic modalities,
together with the underlying resource-based interpretation of
contexts.  We keep the key intuition that ``a resource has the
capability to make demonic/angelic decisions'', and then define a
\emph{contextual capability algebra} as a Kripke resource algebra to
support a bunched type system.  Beyond supporting demonic and angelic
modalities as first-class operators, our logic is also equipped with a
sum operator for sum types, as well as a \emph{dependence modality} to
handle dependencies in refinement types. 

%% \subsection{Store and Contextual Capabilities}

%% As discussed in \autoref{sec:context-typing}, a reified environment
%% $\mathcal{E}(e)$ can provide multiple assignments for free variables
%% in $e$, which enable potential demonic or angelic decisions. For
%% example, environment
%% \[\zlet{x}{\Code{positive\_gen}()}{\boxed{x +  1}}\]
%% provides assignments for variable $x$ in term $x + 1$.  It
%% has the capability to either angelically bind $x$ to be $3$, when used
%% for reachability verification $\vdash x + 1 : \urt{\Nat}{\vnu = 4}$,
%% or to demonically instanitate $x$ with all positive numbers in a safety
%% judgement, $\vdash x + 1 : \ort{\Nat}{\vnu > 0}$.

%% %% On the other
%% %% hand, if an environment has a singleton capability, angelic and
%% %% demonic interpretations collapse into the same
%% %% interpretation. Following this idea, we define this capability as a
%% %% non-empty set of stores that provides assignments for variables.
%% \newcommand{\env}{\rho}

%% \begin{definition}[Environment]\label{def:env}
%%   Assume $\Var$ and $\Val$ range over variable and value domains,
%%   resp. An environment $\rho$ is a partial map from variables to
%%   values, i.e., $\Var \rightharpoonup \Val$.
%%   %% We write $[\overline{x_i
%%   %%     \mapsto v_i}]$ as a literal environment; we write $\env(x)$ for the
%%   %% value of variable $x$ in environment $\rho$.
%%   We define the following standard operations on environments:
%% \begin{itemize}
%% %  \item \emph{Domain}: We write $\Dom{\env}$ for its domain, i.e., the set of variables where the store is defined.
%%   \item \emph{Compatibility}: Write $\mapCompat{\env_1}{\env_2}$ when $\env_1$ and $\sigma_2$ are
%%     \emph{compatible}, i.e.\ they agree on the intersection of their domains.
%%     \[
%%       \mapCompat{\env_1}{\env_2}
%%       \doteq
%%       \forall x \in \Dom{\env_1} \cap \Dom{\env_2}.\; \env_1(x) = \env_2(x).
%%     \]
%% Note that the concatenation of two environment is defined as a map union
%% $\env_1 \cup \env_2$, which is exchangeable when
%% $\mapCompat{\env_1}{\env_2}$ holds.
%%   \item \emph{Projection}: Fix finite $X \subseteq_{\mathrm{fin}}
%%     \Var$ and an environment $\env$; we write $\projectTo{\env}{X}$ for
%%     the projection of $\env$ to $\Dom{\env} \cap X$, i.e.,
%%   \begin{align*}
%%     \projectTo{\sigma}{X} \doteq \{ [x\mapsto \sigma(x)] \mid x \in \Dom{\sigma} \cap X \}.
%%   \end{align*}
%%   \item \emph{Refinement}: The projection operation can derive a
%%     partial order refinement relation. Write $\env_1 \sqsubseteq
%%     \env_2$ when $\env_1$ is a refinement of $\env_2$, i.e.,
%%     $\env_1 = \projectTo{\env_2}{\Dom{\env_1}} $.
%% \end{itemize}
%% \end{definition}

%% The refinement relation $\sqsubseteq$ is important for preserving the
%% semantics of program execution. For example, if a term $e$ can reduce
%% to a value $v$ under environment $\env_1$, then it should also be able to
%% reduce to $v$ under any environment $\env_2$ where $\env_2$ is a
%% refinement of $\env_1$.

\begin{definition}[Contextual capability]\label{def:contextual-capability}
A contextual capability $r$ is a power set of environments,
i.e., $\mathcal{P}(\Code{Var} \rightharpoonup
\Code{Value})$. Such capabilities are well-formed if they are
non-empty and all its environments share the same domain.
%% \begin{align*}
%%   \Code{WellFormed}(r) \doteq r \not= \emptyset \text{ and } \forall \env,\;\env' \in r. \Dom{\env} = \Dom{\env'}
%% \end{align*}\noindent
%% We write $\Res \doteq \{r ~|~ \Code{WellFormed}(r) \}$ for the domain
%% of all well-formed contextual capabilities.
In the rest of the paper,
we assume that contextual capabilities are always well-formed.  
%% Well-formedness also allows us to safely extend environment definitions to
%% contextual capabilities in the obvious way.
%% That is,
%% \begin{itemize}
%%   \item \emph{Domain}: $\Dom{r} \doteq \Dom{\sigma}$ for arbitrary $\sigma \in r$, since $r$ is non-empty and shares the same domain.
%%   \item \emph{Compatibility}: Write $\mapCompat{r_1}{r_2}$ when all stores in $r_1$ are compatible with all stores in $r_2$, i.e.,
%%     \[
%%       \mapCompat{r_1}{r_2}
%%       \doteq
%%       \forall \sigma_1 \in r_1.\;\forall \sigma_2 \in r_2.\; \mapCompat{\sigma_1}{\sigma_2}
%%     \]
%%   \item \emph{Projection}: The projection operation $\projectTo{r}{X}$ is defined piecewise, i.e.,
%%   \begin{align*}
%%     \projectTo{r}{X} \doteq \{ \projectTo{\sigma}{X} \mid \sigma \in r \}.
%%   \end{align*}
%%   \item \emph{Refinement}: the refinement relation $r_1 \sqsubseteq r_2$ holds iff $r_1 = \projectTo{r_2}{\Dom{r_1}} $.
%% \end{itemize}
\end{definition}

\begin{example}[Kripke Semantics]\label{ex-algebraic-semantics} In Kripke semantics, type judgements hold within different
worlds.  In our setting, such worlds are capabilities.  Interpreting
a capability as a denotation of type context, i.e., $r \models
\denot{\Gamma}$, we have:
\begin{align*}
  &r_x \doteq \{[x \mapsto 0],\, [x \mapsto 1]\}
  \\& r_x \models \denot{x {:} \ort{\Nat}{\vnu < 3}} \quad \judgok \quad r_x \models \denot{x {:} \urt{\Nat}{\vnu < 2}}  \quad \judgok
\end{align*}\noindent
Here, $r_x$ is consistent with the first demonic context since all
bindings of $x$ in its different environments are less than $2$; and
is also consistent with the second context since bindings $[x\mapsto
  1]$ and $[x\mapsto 2]$ support angelic decisions consistent with the
qualifier $\vnu < 2$.

The denotation of a typing judgement $\denot{\Gamma \vdash e: \tau}$
is a proposition that holds if the result set of all free variables in
$e$ over all environments in $r$ a contextual capability needs to be
consistent with type $\tau$.'' Let $r_x' \doteq \{[x \mapsto 0],\, [x
  \mapsto 1], [x \mapsto 2]\}$, which means term $x$ can result in a
set of values $\{0, 1, 2\}$ under this capability. Then, we have
\begin{align*}
  r_x \models \denot{x - 1 : \ort{\Nat}{\vnu < 1}} \quad \judgfail \quad r_x \models \denot{x - 1 : \urt{\Nat}{\vnu < 2}}  \quad \judgok
\end{align*}\noindent
The first judgement fails since the result set of $x$ is $\{0, 1\}$,
which is not less than $1$, while the second judgement holds since the
capability allows an angelic choice of all natural numbers less than $2$.
\end{example}

\subsection{A Resource Algebra}

Our context logic is a bunched implication logic (BI-logic) based on a
bunched implication algebra (BI-algebra), which is a partial
commutative monoid with a pre-order that satisfies the bifunctoriality
property~\cite{...}. Moreover, in order to support sum types $\tau_1
\ctoplus \tau_2$, we also add a sum operation on contextual
capabilities.

\begin{definition}[Contextual Capabilities Product]\label{def:contextual-capabilities-product}
  The product of two contextual capabilities $r_1 \times r_2$ is defined
  as the following  partial binary operation:
  \[
    r_1 \times r_2 \;\doteq\;
    \begin{cases}
      \{\, \env_1 \cup \env_2 \mid \env_1 \in r_1,\; \env_2 \in r_2 \,\}
        & \text{if } \mapCompat{r_1}{r_2}, \\
        \text{undefined}
        & \text{otherwise.}
      \end{cases}
  \]
\end{definition}

\begin{definition}[Contextual Capabilities Sum]\label{def:contextual-capabilities-sum}
  The sum of two contextual capabilities $r_1 + r_2$ is defined
  as the following partial binary operation:
  \[
    r_1 + r_2 \;\doteq\;
    \begin{cases}
      r_1 \cup r_2
        & \text{if } \Dom{r_1} = \Dom{r_2}, \\
        \text{undefined}
        & \text{otherwise.}
      \end{cases}
  \]
\end{definition}

Note that these two operations are \emph{partial}: two resources
cannot be multiplied unless they are compatible, and two resources
cannot be summed unless they have the same domain.

\begin{example}[Capability Disjointness]\label{ex-capability-disjointness} The capability product means these two capabilities are disjoint. Consider the following product capability $r_x \times r_y$:
\begin{align*}
  &r_x \doteq \{[x \mapsto 1],\, [x \mapsto 2]\} \quad r_y \doteq \{[y \mapsto 3],\, [y \mapsto 4]\} 
  \\&r_x \times r_y
    = \left\{
      \begin{array}{@{}l@{}}
        [x \mapsto 1,\, y \mapsto 3],\,
        [x \mapsto 1,\, y \mapsto 4],\\[\jot]
        [x \mapsto 2,\, y \mapsto 3],\,
        [x \mapsto 2,\, y \mapsto 4]
      \end{array}
    \right\}.
\end{align*}\noindent
We thus have:
\begin{align*}
  &r_x \times r_y \models \denot{x : \urt{\Nat}{1 \leq \vnu \leq 2}} \quad \judgok \quad r_x \times r_y \models \denot{y : \urt{\Nat}{3 \leq \vnu \leq 4}} \quad \judgok
  \\&r_x \times r_y \models \denot{y - x : \urt{\Nat}{1 \leq \vnu \leq 3}} \quad \judgok
\end{align*}\noindent
where $y - x$ can cover multiple values since $x$ and $y$ are disjoint. However, consider another capability:
\begin{align*}
  &r_{xy} \doteq \{[x \mapsto 1,\, y \mapsto 3],\, [x \mapsto 2,\, y \mapsto 4]\}
  \\&r_{xy} \models \denot{x : \urt{\Nat}{1 \leq \vnu \leq 2}} \quad \judgok \quad r_{xy} \models \denot{y : \urt{\Nat}{3 \leq \vnu \leq 4}} \quad \judgok
  \\&r_{xy} \models \denot{y - x : \urt{\Nat}{1 \leq \vnu \leq 3}} \quad \judgfail
\end{align*}\noindent
Note that term $y - x$ can only reduce to $2$ under $r_{xy}$, thus the last type judgement does not hold. We can also say $x$ and $y$ are not disjoint in $r_{xy}$, i.e., this capability cannot be represented as a product of two capabilities.
\end{example}


\begin{definition}[Contextual Capabilities Algebra]\label{def:contextual-capabilities-algebra}
  The tuple $(\Res,{+},{\times},\mathbf{1})$ forms a \emph{partial
    commutative semiring without zero} where the carrier set $\Res$ is
  the domain of all well-formed contextual capabilities, and the unit
  element $\mathbf{1}$ of product is the singleton capability
  $\{\emptyset\}$. 
%%   \SJ{The following holds by the definition of a semiring, and doesn't need
%%     to be formally written here.}

%%   The following algebraic properties hold:
%%   \begin{itemize}
%%   \item Identity: $\forall r \in \Res. r \times \mathbf{1} = \mathbf{1} \times r = r$.
%%   \item Commutative (when operations are defined): 
%%   \begin{align*}
%%     \forall r_1\;r_2 \in \Res. r_1 + r_2 = r_2 + r_1 \\
%%     \forall r_1\;r_2 \in \Res. r_1 \times r_2 = r_2 \times r_1
%%   \end{align*}
%%   \item Associative (when operations are defined): 
%%   \begin{align*}
%%     \forall r_1\;r_2\;r_3 \in \Res. (r_1 + r_2) + r_3 = r_1 + (r_2 + r_3) \\
%%     \forall r_1\;r_2\;r_3 \in \Res. (r_1 \times r_2) \times r_3 = r_1 \times (r_2 \times r_3)
%%   \end{align*}
%%   \end{itemize}
%% \end{definition}


In order to enable Kripke semantics for bunched implication, this semiring needs to be equipped with a partial order $\leq$ for Kripke's Monotonicity, i.e.,
\begin{align*}
  m_1 \leq m_2 \land m_1 \models \denot{e : \tau} \implies m_2 \models \denot{e : \tau}
\end{align*}\noindent
where $\denot{e : \tau}$ is the denotation of type judgement
(informally). At first glance, the subset relation $\subseteq$ or
superset relation $\supseteq$ over capabilities are natural choices
for the partial order; however, neither can simultaneously model both
angelic and demonic interpretations.
\begin{align*}
  &m_1 \leq m_2 \land m_1 \models \denot{x : \ort{\Nat}{\vnu > 0}} \implies m_2 \models \denot{x : \ort{\Nat}{\vnu > 0}} \\
  &m_1 \leq m_2 \land m_1 \models \denot{x : \urt{\Nat}{\vnu > 0}} \implies m_2 \models \denot{x : \urt{\Nat}{\vnu > 0}}
\end{align*}\noindent
As shown above, since angelic and demonic contexts reflect different
properties (must vs. may), the two judgements above cannot hold at the
same time, regardless of how we interpret $\leq$.  The intuition is
that a Kripke partial order should be neutral, to avoid predefining
either an over- or under-approximation semantics; otherwise we cannot
lift angelic and demonic modalities as first-class operators. Instead,
our algebra defines this partial order in terms of the refinement
relation $\sqsubseteq$ over capabilities.

\begin{definition}[Bifunctoriality]\label{def:bifunctoriality}
  The refinement relation $\sqsubseteq$ over capabilities satisfies the bifunctoriality property, i.e.,
  \begin{align*}
    \forall r_1\;r_2\;r_1'\;r_2' \in \Res. r_1 \sqsubseteq r_1' \land r_2 \sqsubseteq r_2' &\implies r_1 \times r_2 \sqsubseteq r_1' \times r_2' 
    \\& (\text{when } r_1 \times r_2 \text{ and } r_1' \times r_2' \text{ are defined.})
  \end{align*}
\end{definition}

\begin{theorem}[BI-algebra]\label{thm:bi-algebra} The semiring $(\Res,{+},{\times},\mathbf{1})$ with partial order $\sqsubseteq$ forms a BI-algebra.
\end{theorem}

% Intuitively, the bifunctoriality property means that if two pieces of context have fewer capabilities than others, then their product also has fewer capabilities. As a special case, we always have $r_1 \sqsubseteq r_1 \times r_2$ and $r_2 \sqsubseteq r_1 \times r_2$ when $r_1 \times r_2$ is defined. However, the reverse property does not hold in general. That is, for resources with $r_1 \sqsubseteq r_{12}$ and $r_2 \sqsubseteq r_{12}$, we do not generally expect $r_1 \times r_2 \sqsubseteq r_{12}$.

% \begin{example}[Bifunctoriality]\label{ex-bifunctoriality}
% Let $r_x \doteq \{\,[x \mapsto 1],\, [x \mapsto 2]\,\}$, so
% $\Dom{r_x} = \{x\}$.
% Let $r_y \doteq \{\,[y \mapsto 3],\, [y \mapsto 4]\,\}$, with $\Dom{r_y} = \{y\}$.
% Then
% \[
%   r_x \times r_y
%   = \left\{
%     \begin{array}{@{}l@{}}
%       [x \mapsto 1,\, y \mapsto 3],\,
%       [x \mapsto 1,\, y \mapsto 4],\\[\jot]
%       [x \mapsto 2,\, y \mapsto 3],\,
%       [x \mapsto 2,\, y \mapsto 4]
%     \end{array}
%   \right\}.
% \]
% It is easy to see that $r_x \sqsubseteq r_x \times r_y$ and $r_y \sqsubseteq r_x \times r_y$.
% Now take
% $r_{xy} \doteq \{\, [x \mapsto 1,\, y \mapsto 3],\, [x \mapsto 2,\, y \mapsto 4]\,\}$.
% We have
% $\projectTo{r_{xy}}{\{x\}} = r_x$ and $\projectTo{r_{xy}}{\{y\}} = r_y$,
% so by the definition of~$\sqsubseteq$,
% $r_x \sqsubseteq r_{xy}$ and $r_y \sqsubseteq r_{xy}$. However, we cannot refine $r_x \times r_y$ to $r_{xy}$, and vice versa. Actually, we have
% \[
%   r_x \times r_y
%   = r_{xy} + \{\, [x \mapsto 1,\, y \mapsto 4],\, [x \mapsto 2,\, y \mapsto 3] \,\},
% \]
% In general, the product is the one that contains more stores among all capabilities that can be refined by the left and right capabilities of the product.
% \end{example}

\subsection{Context Logic}

\begin{definition}[Context Logic]\label{def:context-logic} The syntax of context logic is defined as follows:
  \begin{align*}
    P, Q &::= \chtrue ~|~ \chfalse ~|~ P \chland P ~|~ P \chlor P ~|~ P \chimpl P ~|~ P \chstar P ~|~ P \chwand P ~|~ \forall x. P ~|~ \exists x. P ~|~ 
    \\& \qquad \Code{Atom}(\phi) ~|~ \demonic P ~|~ \angelic P ~|~ P \choplus P ~|~ \forallM{x} P
  \end{align*}\noindent
The forcing semantics of basic context logic is defined as follows:
\begin{itemize}
\item $r \models \chtrue$ always;
\item $r \models \chfalse$ never;
\item $r \models P \chland Q$ iff $r \models P$ and $r \models Q$;
\item $r \models P \chlor Q$ iff $r \models P$ or $r \models Q$;
\item $r \models P \chimpl Q$ iff for all $r'$, if $r \sqsubseteq r'$ and $r' \models P$ then $r' \models Q$;
\item $r \models P \chstar Q$ iff $\exists r_1\;r_2.\; r_1 \times r_2 \sqsubseteq r$ and
  $r_1 \models P$ and $r_2 \models Q$.
\item $r \models P \chwand Q$ iff $\forall r'.\; \mapCompat{r'}{r}$ and $r' \models P$ implies
  $r' \times r \models Q$.
\item $r \models \forall x.P$ iff $x\notin \Dom{r}$ and $\forall r'. x\in \Dom{r'} \text{ and }r \sqsubseteq r' \text{ implies } r' \models P$.
\item $r \models \exists x.P$ iff $x\notin \Dom{r}$ and $\exists r'. x\in \Dom{r'} \text{ and }r \sqsubseteq r' \text{ and } r' \models P$.
\item $r \models \Code{Atom}(\phi)$ iff $\projectTo{r}{\Code{FV}(\phi)} = \{\sigma ~|~ \Dom{\sigma} = \Code{FV}(\phi) \land \sigma(\phi) \text{ is true} \}$;
\item $r \models \demonic P$ iff $\exists r'.\; r \subseteq r'$ and
$r' \models P$;
\item $r \models \angelic P$ iff $\exists r'.\; r' \subseteq r$ and
$r' \models P$.
\item $r \models P \choplus Q$ iff $\exists r_1\;r_2.\; r_1 + r_2 \sqsubseteq r$ and
  $r_1 \models P$ and $r_2 \models Q$.
\item $r \models \forallM{x}P$ iff $\forall [x\mapsto v_x] \in \projectTo{r}{\{x\}}. \fiberOver{r}{[x\mapsto v_x]} \models P[x\mapsto v_x]$;
\end{itemize}
We still write $p \models q$ iff for all $r \in \Res$, $r \models p$ implies $r \models q$.
\end{definition}

The context logic is a bunched implication logic based on our contextual capabilities algebra. Most logic operators are standard; the rest of this subsection focuses on atomic proposition ($\Code{Atom}$), demonic ($\;\demonic$) and angelic ($\;\angelic$) modalities, sum operator ($\choplus$), as well as reference modality ($\forallM{x}$) that helps express semantics of dependent refinement type system.

\paragraph{Atomic Propositions} The atomic proposition $\Code{Atom}(\phi)$ in our logic is lifted from a standard proposition $\phi$ (e.g., the type qualifier defined in \autoref{sec:context-typing}). For any proposition $\phi$ that is closed under a store, we lift it into a resource that \emph{exactly} contains all stores that satisfy $\phi$. As mentioned before, our logic should be neutral (i.e., not overapproximation or underapproximation) to support demonic and angelic modalities as first-class citizens.

\paragraph{Demonic and Angelic Modality} In order to enable arbitrary combinations of safety and reachability reasoning, we introduce demonic modality $\demonic$ and angelic modality $\angelic$ as operators in our context logic. Similar to previous works~\cite{OHP19}, the demonic modality takes an overapproximation of the capability (upper bound), while the angelic modality takes an underapproximation of the capability (lower bound).

\begin{example}[Approximation modality]
  We have the following facts:
  \begin{align*}
    &\{ [x\mapsto 1]; [x\mapsto 2] \} \models \Code{Atom}(1 \leq x \leq 2)\judgok
    &\{ [x\mapsto 1; y \mapsto 3]; [x\mapsto 2; y \mapsto 4] \} \in \Code{Atom}(1 \leq x \leq 2)\judgok
  \end{align*}\noindent
  Here we require the capability exactly has stores that make $x$ be $1$ and $2$. Although the capability is allowed to provide information about other variables like $y$, the capability cannot lose any case or include extra cases:
  \begin{align*}
    &\{ [x\mapsto 1] \} \models \Code{Atom}(1 \leq x \leq 2)\judgfail
   &\{ [x\mapsto 1]; [x\mapsto 2]; [x\mapsto 3] \} \models \Code{Atom}(1 \leq x \leq 2)\judgfail
  \end{align*}\noindent
  On the other hand, the demonic and angelic modality can enable the approximation.
  \begin{align*}
    &\{ [x\mapsto 1] \} \models \demonic \Code{Atom}(1 \leq x \leq 2)\judgok
   &\{ [x\mapsto 1]; [x\mapsto 2]; [x\mapsto 3] \} \models \angelic \Code{Atom}(1 \leq x \leq 2)\judgok
  \end{align*}\noindent
  where these two formulas are the denotation of types $\ort{\Nat}{1 \leq x \leq 2}$ and $\urt{\Nat}{1 \leq x \leq 2}$.
  We also have the following facts for the two capabilities in Example~\ref{ex-capability-disjointness}: 
  \begin{align*}
    &r_x \times r_y \models (\angelic \Code{Atom}(1 \leq x \leq 2)) \chstar (\angelic \Code{Atom}(3 \leq y \leq 4))  \quad \judgok
   \\&r_{xy} \models (\angelic \Code{Atom}(1 \leq x \leq 2)) \chland (\angelic \Code{Atom}(3 \leq y \leq 4))  \quad \judgok
  \end{align*}\noindent 
  The $r_x \times r_y$ can model the second formula; however, $r_{xy}$ cannot model the first formula that requires disjointness.
\end{example}


% \begin{example}[Safety and Reachability Verification] We still consider capability $r_{xy}$ in Example~\ref{ex-basic-context-logic}; we have:
%   \begin{align*}
%     &r_{xy} \models (1 \leq x \leq 2)
%     && r_{xy} \not\models \demonic (x = 1)
%     && r_{xy} \models \demonic (0 \leq x \leq 4)
%     \\&r_{xy} \models (1 \leq x \leq 2)
%     && r_{xy} \models \angelic (x = 1)
%     && r_{xy} \not\models \angelic (0 \leq x \leq 4)
%   \end{align*}
% In the first line, $r_{xy}$ is allowed to satisfy qualifiers that are overapproximations of the values it can actually reach. This is safety verification: the assertion $\demonic (x = 1)$ fails since a demonic context with capability $r_{xy}$ can adversarially assign $x$ to be $2$, which is unsafe. On the other hand, assertion $\angelic (x = 1)$ holds since there exists an store in $r_{xy}$ that makes $x = 1$ hold; assertion $\angelic (0 \leq x \leq 4)$ fails since $x = 0$ is unreachable from $r_{xy}$.
% \end{example}

% Reader may be curious about why standard refinement type system doesn't consider the disjointness over context peices. In fact, the $P \chland Q$ and $P \chstar Q$ is collapsed into the same meaning when both $P$ and $Q$ are demonic formula. For example
% \begin{align*}
%   \demonic (1 \leq x \leq 2) \chland \demonic (3 \leq y \leq 4)  &\equiv \demonic (1 \leq x \leq 2 \land 3 \leq y \leq 4)
% \\&\equiv \demonic (1 \leq x \leq 2) \chstar \demonic (3 \leq y \leq 4)
% \end{align*}\noindent
% Actually, the demonic modalit takes overapproximation semantics, which leads 

\paragraph{Sum operator} The $P \choplus Q$ operation is defined as the sum operation on capabilities, which can be interpreted as ``if the capability is divided into two parts, one part satisfies $P$ and the other part satisfies $Q$.'' 

\begin{example}[Sum operator] The examples in \autoref{sec:context-typing} allow us to divide current type context $x{:}\urt{\Nat}{\top}$ to fit condition expression $x == 0$ of the program, which requires:
\begin{align*}
  &\{[x\mapsto 0], [x\mapsto 1], ... \} = \{[x\mapsto 0]\} + \{[x\mapsto 1], ... \}
  \\&\{[x\mapsto 0]\} \models \angelic \Code{Atom}(x = 0) \judgok
  \\&\underline{\{[x\mapsto 1], ... \} \models \angelic \Code{Atom}(x > 0) \judgok}
  \\&\{[x\mapsto 0], [x\mapsto 1], ... \} \models \angelic \Code{Atom}(x = 0) \choplus \angelic \Code{Atom}(x > 0) \judgok
\end{align*}\noindent
Note that our well-formed capability (i.e., $\Res$) is nonempty set by definition, which avoids the vacously proof in reachability verification. Otherwise, consider the control flow example in \autoref{sec:context-typing}:
\begin{align*}
  \denot{x{:}\urt{\Nat}{\vnu > 0}} \ctsubseteq{\{x\}} \denot{x{:}\ort{\Nat}{\bot} \ctoplus x{:}\urt{\Nat}{\vnu > 0}}  \quad \text{(without nonemptyness constraint)}
\end{align*}\noindent
where $x{:}\ort{\Nat}{\bot}$ is a type context that can be satisified by a empty capability (i.e., emptyset of stores). This ill sum context means that one typecheck an unreachable control flow with a bottom type (e.g., $x{:}\urt{\Nat}{\vnu > 0}$) where any return type is then allowed.
\end{example}

\paragraph{Variable dependency} Although a neutral atomic proposition allows us to lift demonic and angelic modalities as first-class citizens, it creates additional difficulty in expressing variable dependency. Consider a program $\zlet{y}{x + 1}{...}$ where $y$ always has assignment $x + 1$. If we directly encode this relation as a logical formula, we have:
\begin{align*}
  r \models \Code{Atom}(y = x + 1)  \iff \projectTo{r}{\{x, y\}} = \{ [x \mapsto 0; y \mapsto 1], [x \mapsto 1; y \mapsto 2], ... \}
\end{align*}\noindent
Since $\Code{Atom}$ has exact semantics, capability $r$ must contain the case where $x$ is any natural number. Then, the following fact does not hold:
\begin{align*}
  \Code{Atom}(1 \leq x \leq 2) \chland \Code{Atom}(y = x + 1) &\not\equiv \Code{Atom}(1 \leq x \leq 2 \land y = x + 1) 
\end{align*}\noindent
Actually, the left formula is equivalent to $\chfalse$, since $\Code{Atom}(1 \leq x \leq 2)$ requires ``$x$ must be $1$ or $2$'' while $\Code{Atom}(y = x + 1)$ requires $x$ to range over all natural numbers.

In short, the neutral interpretation of atomic predicate means that $\Code{Atom}(y = x + 1)$ cannot capture the constraint about $x$ from other conjunctive clauses. Adding an explicit demonic modality can simulate the behavior as expected, i.e.,
\begin{align*}
  \demonic (1 \leq x \leq 2) \chland \demonic (y = x + 1) \equiv \demonic (1 \leq x \leq 2 \land y = x + 1) 
\end{align*}\noindent
However, it does not work for angelic modality. Intuitively, we should tell pure qualifier $y = x + 1$ that $x$ should depend on the actual capability of $x$ (e.g., $x$ is $1$ or $2$), which can be defined by an explicit operator in our logic.

\begin{definition}[Fibers]\label{def:fibers}
  For a finite variable set $X \subseteq_{\Code{fin}} \Var$, consider a store $\sigma \in \projectTo{r}{X}$, which is one result of projection. We write $\fiberOver{r}{\sigma}$ for the corresponding fiber, i.e.,
  \begin{align*}
    \fiberOver{r}{\sigma} \doteq \{\, \sigma' \mid \sigma' \in r \text{ and } \projectTo{\sigma'}{\Dom{\sigma}} = \sigma \,\}.
  \end{align*}\noindent
  In short, the fiber $\fiberOver{r}{\sigma}$ is a subset of $r$ whose projection to variable set $X$ equals $\sigma$.
\end{definition}

\begin{definition}[Reference Operator]\label{def:context-logic-reference-operator} We define a reference operator $\forallM{x}{P}$, which means $x$ is a referred variable on the current capability, i.e.,
  \begin{itemize}
    \item $r \models \forallM{x}P$ iff $\forall [x\mapsto v_x] \in \projectTo{r}{\{x\}}. \fiberOver{r}{[x\mapsto v_x]} \models P[x\mapsto v_x]$;
  \end{itemize}
The notation $\forallM{\overline{x_i}} P$ is shorthand for
  $\forallM{x_1}\,\forallM{x_2}\,\cdots\,\forallM{x_n}\,P$.
\end{definition}

The intuition of reference operator $\forallM{x}{P}$ is quite simple: we first select all allowed assignments of $x$ in the current capability (i.e., $v_x$), and then check $P$ under the corresponding fiber. Importantly, we substitute variable $x$ with concrete assignment $v_x$ in logical formula $P$ so that pure qualifiers (e.g., from $y = x + 1$ to $y = v_x + 1$) do not need to consider all possible assignments of $x$.

\begin{example}[Variable Reference] With the help of the reference operator, we can express variable dependency across clauses. Consider the following logic formula, which is the type denotation of an angelic type with a referred variable $x$:
\begin{align*}
  \denot{y{:}\urt{\Nat}{0 \leq \vnu < x}} = \forallM{x} \angelic (0 \leq y < x)
\end{align*}\noindent
It means that the variable $y$ must reach all values in the range $[0, x)$ whatever $x$ is, then we have:
  \begin{align*}
    &r_1 \doteq \{[x \mapsto 1; y \mapsto 0]\} \quad r_2 \doteq \{[x \mapsto 2; y \mapsto 0]; [x \mapsto 2; y \mapsto 1] \}\\
    &r_1 \models \forallM{x} \angelic (0 \leq y < x) \judgok \quad r_1 \models \forallM{x} \angelic (0 \leq y < x)\judgok   \quad r_1 + r_2 \models \forallM{x} \angelic (0 \leq y < x)  \judgok 
  \end{align*}
Furthermore, we have
\begin{align*}
  \denot{x{:}\urt{\Nat}{1 \leq \vnu \leq 2} \ctcomma y{:}\urt{\Nat}{0 \leq \vnu < x}} =
  \angelic (1 \leq x \leq 2) \chland \forallM{x} \angelic (0 \leq y < x) 
\end{align*}\noindent
where both $x$ and $y$ have angelic interpretation while $y$ is dependent on $x$. Note that in dependent type system, $\Gamma_1 \ctcomma \Gamma_2$ is not exchangeable since only $\Gamma_2$ can refer to bindings in $\Gamma_1$; however, logical conjunction $\chland$ does not have this restriction. This dependency in refinement type system is encoded into reference operator $\forallM{x}$, which makes $\forallM{x} \angelic (0 \leq y < x) \chland \angelic (1 \leq x \leq 2)$ also a valid denotation.
\end{example}

% \begin{example}[Simulating Hoare-logic and Incorrectness-logic] With the help of demonic and angelic modalities, we can simulate Hoare-logic-style safety verification and Incorrectness-style reachability verification~\cite{OHP19}. Consider a command $\Code{y := x + 1}$ that reads the value of $\Code{x}$ and assigns $\Code{y}$ as $\Code{x} + 1$, which can be expressed as a logical formula $\forallM{x} (y = x + 1)$. Then a Hoare triple $\{x > 0\}\;\Code{y := x + 1}\;\{y > 0\}$ says ``if $x$ is greater than $0$, then after executing $\Code{y := x + 1}$, $y$ is greater than $0$'' and can be expressed as:
%   \begin{align*}
%     \demonic (x > 0) \models \forallM{x} (y = x + 1) \chimpl \demonic (y > 0)
%   \end{align*}\noindent
% where the semantics of the command is encoded as a precise formula, while pre- and post-conditions are encoded as demonic formulas.
% On the other hand, an Incorrectness triple $[x > 0]\;\Code{y := x + 1}\;[y > 2]$ says ``for all $y$ that is greater than $2$, there exists an $x$ that is greater than $0$ and $x + 1$ can reach such $y$'' can be expressed as:
%   \begin{align*}
%     \angelic (x > 0) \models \forallM{x} (y = x + 1) \chimpl \angelic (y > 2)
%   \end{align*}\noindent
% \end{example}

\begin{theorem}[Kripke's Monotonicity] 
  \begin{align*}
    r \sqsubseteq r' \text{ and } r \models P \text{ implies } r' \models P
  \end{align*}\noindent
\end{theorem}

\subsection{Type Denotation}

Following the intuition introduced in the previous subsection, we provide a formal definition of type denotation. We define $\denot{\tau}_\Sigma: \Code{Term} \sarr \Code{Formula}$ under a basic type context $\Sigma$, where $r \models \denot{\tau}_\Sigma\;e$ means that $e$ is an inhabitant of type $\tau$ under capability $r$. Like other refinement type systems (e.g., Liquid Types), we also require terms to be total (no divergence) and well-typed in basic type system $\Sigma$. All of these side conditions are encoded in our logic:
\begin{align*}
  P_{\Code{guard}} \doteq \forallM{\Dom{\Sigma}} \Code{Atom}(\Code{Total}(e) \land \Sigma \basicvdash e : \eraserf{\tau} \land \Sigma \wfvdash \tau)
\end{align*}\noindent
Note that we put variable reference $\forallM{\Dom{\Sigma}}$ before the atomic proposition; otherwise, $\Code{Atom}(x{:}\Nat \basicvdash x : \Nat)$ will require capability to contain all possible assignments of $x$ (i.e., all natural numbers), which makes this basic type judgement hold. In the following definition of type denotation, all cases contain this guard, which is omitted for brevity.

\begin{definition}[Term Denotation] The denotation of term $e$ is based on its operational semantics, where a fresh variable $x$ stores all values that $e$ can reduce to under all stores in current capability.
  \begin{align*}
    \denot{\mstep{e}{x}}_\Sigma &\doteq \forallM{\Dom{\Sigma}} \Code{Atom}(\mstep{e}{x})
  \end{align*}\noindent
This denotation has the same reason to have a prefix reference modality. It is important to ensure the referred variable cannot contain fresh variable $x$, especially for a nondeterministic language where $e$ may reduce to multiple values. For example, for a boolean generator function $\Code{boolGen()}$ under empty context, a wrong denotation would be:
\begin{align*}
  \forallM{x} \Code{Atom}(\mstep{\Code{boolGen()}}{x})
\end{align*}\noindent
This denotation requires a resource that is single-valued at $x$ to model $\Code{Atom}(\mstep{\Code{boolGen()}}{x})$, which has exact semantics, and this is impossible.
\end{definition}

\begin{definition}[Type Denotation] The type denotation is defined as follows:
\begin{align*}
  \denot{\ort{b}{\phi}}_\Sigma\;e &\doteq \denot{\mstep{e}{\vnu}}_\Sigma \chimpl \forallM{\Code{FV}(\ort{b}{\phi})}\, \demonic \phi \quad \text{($\vnu$ is a fresh variable)}
  \\\denot{\urt{b}{\phi}}_\Sigma\;e &\doteq \denot{\mstep{e}{\vnu}}_\Sigma \chimpl \forallM{\Code{FV}(\urt{b}{\phi})}\, \angelic \phi \quad \text{($\vnu$ is a fresh variable)}
  \\\denot{x{:}\tau_x \ctsarr \tau}_\Sigma\;e &\doteq \denot{\tau_x}_\Sigma\;x \chimpl \denot{\tau}_{\Sigma,x{:}\eraserf{\tau_x}}\;(e\;x) \quad \text{($x$ is a fresh variable)}
  \\\denot{x{:}\tau_x \ctwand \tau}_\Sigma\;x &\doteq \denot{\tau_x}_\Sigma\;x \chwand \denot{\tau}_{\Sigma,x{:}\eraserf{\tau_x}}\;(e\;x) \quad \text{($x$ is a fresh variable)}
  \\\denot{\tau_1 \ctinter \tau_2}_\Sigma\; e &\doteq \denot{\tau_1}_\Sigma\; e \chland \denot{\tau_2}_\Sigma\; e
  \\\denot{\tau_1 \ctunion \tau_2}_\Sigma\; e &\doteq \denot{\tau_1}_\Sigma\; e \chlor \denot{\tau_2}_\Sigma\; e
  \\\denot{\tau_1 \ctoplus \tau_2}_\Sigma\; e &\doteq \denot{\tau_1}_\Sigma\; e \choplus \denot{\tau_2}_\Sigma\; e
  % \\\denot{\tau_1 \ctoplus \tau_2}\; e &\doteq \forall x. \denot{e}_x \chimpl \denot{\tau_1}\;x \choplus \denot{\tau_2}\;x
  % \\\denot{\ctpers \tau}\; e &\doteq \chpers (\denot{\tau}\;e)
\end{align*}\noindent
Denotation of demonic and angelic types is directly encoded into our logic, where we still use reference operator to make type qualifiers aware of referred variables (i.e., $\Code{FV}(\ort{b}{\phi})$ and $\Code{FV}(\urt{b}{\phi})$). Intuitively, it means that ``if we store result of $e$ at fresh variable $\vnu$, then denotation of type qualifier holds with corresponding angelic or demonic modality.'' Function type denotation is also standard: we assume parameter $x$ stores values consistent with parameter type $\tau_x$, then function application $e\;x$ is an inhabitant of return type $\tau$. We do not need substitution, since all variable-reference relations are encoded in the logic formula.
\end{definition}

\begin{definition}{Type context denotation}
  We define $\denot{\Gamma}_\Sigma: \Code{Formula}$: 
  \begin{align*}
    \denot{\emptyset}_\Sigma &\doteq \chtrue
    \\\denot{x{:}\tau}_\Sigma &\doteq \denot{\tau}_\Sigma\;x
    \\\denot{\Gamma_1 \ctcomma \Gamma_2}_\Sigma &\doteq \denot{\Gamma_1}_\Sigma \chland \denot{\Gamma_2}_{(\Sigma, \eraserf{\Gamma_1})}
    \\\denot{\Gamma_1 \ctstar \Gamma_2}_\Sigma &\doteq \denot{\Gamma_1}_\Sigma \chstar \denot{\Gamma_2}_\Sigma
    \\\denot{\Gamma_1 \ctoplus \Gamma_2}_\Sigma &\doteq \denot{\Gamma_1}_\Sigma \choplus \denot{\Gamma_2}_\Sigma
  \end{align*}\noindent
\end{definition}

We will omit the basic type context $\Sigma$ when it is clear from the context.

\begin{example}[Type Denotation] Consider the following type denotation:
  \begin{align*}
    &\denot{x{:}\urt{\Nat}{\vnu > 0}\ctsarr y{:}\urt{\Nat}{\vnu \geq 0} \ctwand \urt{\Nat}{\vnu > x}}\;e
\\=\;& \denot{\urt{\Nat}{\vnu > 0}}\;x \chimpl \denot{y{:}\urt{\Nat}{\vnu \geq 0} \ctwand \urt{\Nat}{\vnu > x}}\;(e\;x)
\\=\;& \denot{\urt{\Nat}{\vnu > 0}}\;x \chimpl \denot{\urt{\Nat}{\vnu \geq 0}}\;y \chwand \denot{\urt{\Nat}{\vnu > x}}\;(e\;x\;y)
\\=\;& (\denot{\mstep{x}{\vnu_1}} \chimpl \demonic \Code{Atom}(\vnu_1 > 0)) \chimpl
\\&\quad (\denot{\mstep{y}{\vnu_2}} \chimpl \angelic \Code{Atom}(\vnu_2 \geq 0))
  \chwand (\denot{\mstep{e\;x\;y}{\vnu_3}} \chimpl \forallM{x}\angelic \Code{Atom}(\vnu_3 > x))
  \end{align*}\noindent
We rewrite the fresh variable $\vnu_1$ and $\vnu_2$ as $x$ and $y$ respectively, which can simplify the formula.
\begin{align*}
=\;& (\chtrue \chimpl \demonic \Code{Atom}(x > 0)) \chimpl (\chtrue \chimpl \angelic \Code{Atom}(y \geq 0))
\chwand (\denot{\mstep{e\;x\;y}{\vnu_3}} \chimpl \forallM{x}\angelic \Code{Atom}(\vnu_3 > x))
\\=\;& \demonic \Code{Atom}(x > 0) \chimpl \angelic \Code{Atom}(y \geq 0)
\chwand (\denot{\mstep{e\;x\;y}{\vnu_3}} \chimpl \forallM{x}\angelic \Code{Atom}(\vnu_3 > x))
\end{align*}\noindent
Then the term $\zzlam{x}{\zzlam{y}{x + y}}$ can have this type under empty context, which requires that
\begin{align*}
  \demonic \Code{Atom}(x > 0) \chstar \angelic \Code{Atom}(y \geq 0) \chland \denot{\mstep{x+y}{\vnu_3}}
  \models \forallM{x}\angelic \Code{Atom}(\vnu_3 > x)
\end{align*}\noindent
This formula is correct, since whatever which number $\vnu_3$ that required to reach by the right hand side and what demonic assignment $x$ is, the context can angelic assign $y$ to $\vnu_3 - x$ to make $x + y$ reach $\vnu_3$. Moreover, since $\vnu_3 > x$, then the value of $y$ (i.e., $\vnu_3 - x$) must be non-negative which is allowed by $\angelic \Code{Atom}(y \geq 0)$. 
\end{example}


\begin{figure}[h!]
  {\small
  {\normalsize \begin{flalign*}
    &\text{\textbf{Semantic subtyping rules}} && \fbox{$\Gamma \vdash \tau \subty \tau$ \quad $\Gamma \ctsubseteq{X} \Gamma$ }
  \end{flalign*}}
\mprooftr{48mm}
{$\forall e. \denot{\Gamma} \models \denot{\tau_1}\;e \chimpl \denot{\tau_2}\;e$}
{Sub}
{$\Gamma \vdash \tau_1 <: \tau_2$}
\quad
\mprooftr{66mm}
{$\forall r. r \models \denot{\Gamma_1} \impl \exists r'. \restrictTo{r}{X} = \restrictTo{r'}{X} \land r' \models \denot{\Gamma_2}$ }
{CtxSub}
{$\Gamma_1 \ctsubseteq{X} \Gamma_2$}
	  }
  \caption{Semantic subtyping relation}
  \label{fig:semantic-subtyping-relation}
  \vspace*{-.1in}
\end{figure}

\paragraph{Semantical subtyping relation} The definition of semantic subtyping relation shown in \autoref{sec:context-typing} shown in \autoref{fig:semantic-subtyping-relation}. The type subtyping reulation are directly encoded into an implication formula in our logic, while the context subtyping relation requires an explicit projection operation $\restrictTo{r}{X} = \restrictTo{r'}{X}$ which indicates the two capabilities are consistent only on the variable set $X$.

\begin{theorem}[Well-formed operator context] Operator context $\Phi$ is well-formed iff for all primitive operators $\primop$ we have:
  {\begin{align*}
      &\Phi(\primop) = \overline{x{:}\ort{b_x}{\phi_x}}\ctsarr \ort{b}{\phi} \ctinter \urt{b}{\phi} \text{ and } \models \denot{\overline{x{:}\ort{b_x}{\phi_x}}} \chiff \denot{\ort{b}{\phi} \ctinter \urt{b}{\phi}}\; (\primop\;\overline{x})
    \end{align*}}\noindent
  \end{theorem}
  
  \begin{theorem}[Fundamental theorem]$\Gamma \ctvdash e : \tau \text{ implies } \denot{\Gamma}_\emptyset \chvdash \denot{\tau}_{\eraserf{\Gamma}}\;e$
  \end{theorem}
  
  \begin{theorem}[Denotational Soundness]
    \begin{align*}
      \emptyset \ctvdash e : \tau \text{ implies for any fresh variable $x$, } \{[x\mapsto v] ~|~ \mstep{e}{v}\} \models \denot{\tau}\; x
    \end{align*}\noindent
  \end{theorem}
